% !TeX root = ../../Book1.tex
% 纲要标题：### 9.4 整体公式的粘合证明
% 纲要位置：几何学计划纲要.md，第 486 行。
% 正文状态：初稿；尚未正式审核。

% 纲要内容开始
% 
% * 适应边界的曲面分解
% * 公共边项抵消及顶点贡献
% * Euler 示性数的出现
% * 闭曲面与带边界曲面的统一公式
% 
% 纲要内容结束

\section{整体公式的粘合证明}
\label{g1:sec:0904}

一个圆盘的局部公式右端是 \(2\pi\)，若把整张曲面切成 \(F\) 个圆盘，直接相加便得到 \(2\pi F\)。但切割也制造了原来没有的边和角：公共边的曲率会抵消，新增角的贡献却必须逐点数清。最后留下的整数并不是圆盘数，而是第六章的 \(V-E+F\)。这一过程还需要先解决一个正则性问题：第六章给出的是拓扑剖分，拓扑边没有可供积分的曲率，不能原样代入本章公式。

\subsection{适应边界的曲面分解}
\label{g1:subsec:0904:01}

下面允许 \(M\) 有分片 \(C^2\) 边界，角点有限，每个角点的内角属于 \((0,2\pi)\)。这里曲面和度量在各局部图中光滑；边界和角点表示在这样的光滑局部环境内截取一个区域，因而内角、边的切向量和协变加速度都有意义。内角为 \(\pi\) 的分段点也允许保留。

\begin{lemma}[坐标圆盘分解]\label{g1:lem:smooth-coordinate-disk-decomposition}
	设 \(M\) 是紧、Hausdorff、第二可数的光滑度量曲面，允许有分片 \(C^2\) 边界及有限个内角属于 \((0,2\pi)\) 的角点。则存在有限个闭圆盘 \(D_1,\ldots,D_F\)，满足：各 \(D_j\) 位于一张光滑坐标图中，边界由有限条 \(C^2\) 正则弧组成，各角的内角属于 \((0,2\pi)\)；圆盘的内部两两不交，合起来覆盖 \(M\)；两个圆盘的公共部分是若干条完整的边及顶点；\(\partial M\) 由分解的边和顶点组成，原有角点均为顶点。
\end{lemma}
\begin{proof}
	我们沿用\cref{g1:lem:finite-network-disk-cells}的构造思路：先造有限曲线网，再取顶点盘、边带，最后把平面区域中的孔连接起来。需要补充的是，所有新边都可以选成分片光滑的，并且每一块仍留在一张坐标图中。

	先在各点附近选带余量的小坐标片。在内点处用圆盘，在光滑边界点处用半圆盘；角点附近，在环境坐标中用小圆周截取两条边界弧围成的扇形区域。后一构造确实给出闭圆盘：两条正则弧从角点出发时，沿各弧的坐标半径在足够短的一段上严格增加，故每条小圆周恰与每侧相交一次；两侧的不同切向射线保证区域有正角。把最小各片的相对内部取成有限覆盖，再为每片保留稍大的同类片，其半径可以在一个开区间内变化而仍包含原来最小的片。这样的余量与\cref{g1:lem:nested-surface-charts}中的三层坐标片起同样作用。

	逐个选择这些较大片的半径，使新增的圆周弧与已经选定的边及 \(\partial M\) 有限次横交，并避开已有顶点。这个选择只需要如下一个实变量事实。若 \(f\in C^1([a,b])\)，则临界值集 \(f(\{f'=0\})\) 可以用总长度任意小的开区间覆盖。事实上，对任意 \(\varepsilon>0\)，在包含 \(\{f'=0\}\) 的开集中令 \(|f'|<\varepsilon\)；把该开集写成不交区间之并，平均值定理说明各像区间的长度不超过相应区间长度的 \(\varepsilon\) 倍，总和不超过 \(\varepsilon(b-a)\)。端点另用任意短的区间覆盖。因此有限个这类临界值集与有限个指定值，不能填满一个开区间。

	对每条已有光滑弧，令 \(f\) 为沿弧的坐标半径平方。只在所选的闭坐标环带上使用此函数；可以把径向函数在更大圆盘外取为常数，光滑延拓后再复合弧的参数，所以弧反复进出坐标图也不会妨碍上述论证。选择一个同时不是各弧临界值、也不是已有顶点半径平方的值。相交处沿弧的半径导数非零，正是两条曲线的切线不同；交点因而孤立。交点集又是紧集，故有限。避开已有顶点也排除了三条新旧边在同一点相交的情形。初始曲线网取为 \(\partial M\)，并把原有角点标出；最后给没有交点的闭曲线添上几个顶点，就得到一个有限的分片光滑嵌入图 \(G\)。

	每个 \(M\setminus G\) 的连通分支都留在某个已选坐标片内。确实，选分支中一点及包含它的坐标片内部；该片的相对边界已在 \(G\) 中，所以在不碰 \(G\) 的连通分支内，不能从片内走到片外。它的闭包也位于这张坐标图中，因为所选闭片与图的边缘之间仍有余量。不过补分支未必是盘，同一条边还可能在其边界行走中出现两次；下面的加宽步骤就是为了消除这种现象。

	在各顶点取两两不交的小坐标盘；边界顶点取相应半盘或角区域。半径取得足够小，每条关联边只以一个短弧进入顶点盘，并横交它的圆弧边界。把余下的边进一步有限分段，使每个闭段在一张坐标图中都是一条函数图像 \(v=f(u)\)；这是正则性与紧性给出的有限选择。在这样的图中，\(|v-f(u)|\le\delta\) 给出一条细带，位于 \(\partial M\) 上的边只取曲面内侧的半条带。可以令带的两端接到相邻顶点盘的边界弧上，并在接点附近调整侧边使其横交圆弧。有限条互不相交的截短边及其他不关联部分都是紧集，逐次缩小带宽即可使不同带的内部不交，且每个顶点盘、每条带都包含在一张坐标图中。顶点盘和带的并记为 \(N\)。它是 \(G\) 的闭邻域，包含 \(\partial M\) 的相对邻域，边界由有限条分片光滑 Jordan 曲线组成。所有接点都采用横交弧，故不会出现零角尖点。

	余下的每个闭区域 \(R\subset M\setminus\operatorname{int}_M N\) 位于原来一个补分支的坐标图中，其边界为有限条互不相交的分片光滑圈。区域只有有限多个：每个非空区域都接到 \(\partial N\) 的某条边界圈，而一条圈的外侧只邻接一个区域。\cref{g1:lem:finite-network-disk-cells}的平面部分说明，每个连通的 \(R\) 是一个外圈内去掉有限个内圈所围闭盘后的区域。若有内圈，在其一个非顶点处与外圈的非顶点处之间画一条内部位于 \(R^\circ\) 的简单分片光滑弧。具体地，先在两端取横交边界的小线段，再用 \(R^\circ\) 中的有限折线连接两段内端，删去折线回路，并在不碰边界的小圆盘内磨圆折线拐角。取这条弧的一条窄闭带，使其短边沿两条边界弧贴合。

	将窄带作为一个独立盘形块取出，剩下区域的两个边界圈合为一个，内圈数减少一；此处保留带的两个不同侧边，不把它们压回同一条割线。故余下区域始终具有真正的、互不相交的 Jordan 边界。有限次以后，每块余下区域只有一个边界圈，由\cref{g1:thm:jordan-schoenflies}是闭圆盘。它们和取出的窄带、顶点盘、原有边带一起给出有限闭盘分解，新增块都仍在所操作的原坐标图中。最后把所有接点统一作为顶点，边界上的分点也同时加入相邻两块。这样公共部分由完整边和顶点组成，原有边界和角点均保留下来。所有边由上述光滑弧构成，角可以选成严格介于 \(0\) 和 \(2\pi\) 之间，得到所需分解。
\end{proof}

这里没有要求分解的边是测地线。普通光滑边已经足够，因为它们的曲率积分会在粘合时抵消；若强求测地三角剖分，反而会额外引入尚未需要的存在性问题。闭盘也不必是三角形，下面按边和角的实际出现次数计数即可。

\subsection{公共边项抵消及顶点贡献}
\label{g1:subsec:0904:02}

先假设 \(M\) 可定向并选定定向，每块 \(D_j\) 的边界都按区域在左的方向行进。在一条内部公共边上，两块给出的行进方向相反。同一个曲线加速度在反向参数化后不变，而左侧的单位切向法向变号，因此两侧测地曲率互为相反数。两次积分使用相同的弧长测度，恰好抵消。

\ProseParagraphBreak

这个抵消其实不需要整张曲面可定向。对每块圆盘分别选一个局部定向，令其边界仍把本块放在左边。写 \(\nu_j\) 为沿边、垂直于边且指向 \(D_j\) 内侧的单位切向量，则正向行进时 \(\nu_j=JT\)，因而
\[
	k_{g,j}=\langle\nabla_TT,\nu_j\rangle.
\]
反向行进将 \(T\) 换成 \(-T\)，但 \(\nabla_{-T}(-T)=\nabla_TT\)；只要固定内侧，右端便不变。公共边两侧的 \(\nu_j\) 相反，所以无论各面选择的定向能否拼成整体定向，两项都抵消。沿原有边界，圆盘内侧与 \(M\) 内侧一致，留下的恰是 \(M\) 的边界曲率积分。于是
\begin{equation}\label{g1:eq:gb-edge-cancellation}
	\sum_{j=1}^F\int_{\partial D_j}k_{g,j}\,\mathrm{d}s
	=\int_{\partial M}k_g\,\mathrm{d}s.
\end{equation}
右端沿各条光滑边求和，不把角点另算成曲率函数的值。

\ProseParagraphBreak

再看角。用 \(V_{\mathrm i},V_{\mathrm b}\) 分别表示内部与边界顶点数，用 \(E_{\mathrm i},E_{\mathrm b}\) 表示内部与边界边数。每次一个圆盘在某顶点占据一个角，都算作一次角的出现，记其内角为 \(\beta_{j,v}\)。在内点，所有小角在切平面内围满一周，总和为 \(2\pi\)；在一个光滑边界点，它们铺满内侧半平面，总和为 \(\pi\)；在原有角点 \(v\)，总和等于 \(M\) 的内角 \(\beta_v\)，而不是 \(\pi\)。\cref{g1:fig:gb-vertex-angles}画的是顶点处带内积的切平面，各扇区就是这些角的切向模型。

% !TeX root = ../Book1.tex
\begin{figure}[htbp]
	\centering
	\begin{tikzpicture}[scale=0.88,font=\small]
		\begin{scope}
			\fill[GeometryBlueLight] (0,0) circle (1.15);
			\draw[GeometryBlue] (0,0) circle (1.15);
			\foreach \a in {0,55,135,210,285}{\draw[GeometryBlue] (0,0)--(\a:1.15);}
			\draw[GeometryBlue] (0.38,0) arc (0:55:0.38);
			\node at (27:0.7) {\(\beta_1\)};
			\fill (0,0) circle (1.3pt);
			\node[align=center] at (0,-1.65) {内点\\\(\sum\beta_j=2\pi\)};
		\end{scope}
		\begin{scope}[xshift=4cm]
			\fill[GeometryBlueLight] (-1.15,0)--(1.15,0) arc (0:180:1.15)--cycle;
			\draw[GeometryBlue] (-1.15,0)--(1.15,0) arc (0:180:1.15);
			\foreach \a in {45,110}{\draw[GeometryBlue] (0,0)--(\a:1.15);}
			\fill (0,0) circle (1.3pt);
			\node at (0,-0.45) {\(\partial M\)};
			\node[align=center] at (0,-1.65) {光滑边界点\\\(\sum\beta_j=\pi\)};
		\end{scope}
		\begin{scope}[xshift=8cm]
			\fill[GeometryBlueLight] (0,0)--(1.15,0) arc (0:125:1.15)--cycle;
			\draw[GeometryBlue] (0,0)--(1.15,0) arc (0:125:1.15)--cycle;
			\foreach \a in {40,83}{\draw[GeometryBlue] (0,0)--(\a:1.15);}
			\draw[GeometryBlue] (0.42,0) arc (0:125:0.42);
			\node at (63:0.72) {\(\beta_v\)};
			\fill (0,0) circle (1.3pt);
			\node[align=center] at (0.1,-1.65) {原有角点\\\(\sum\beta_j=\beta_v\)};
		\end{scope}
	\end{tikzpicture}
	\caption[粘合时的顶点角计数]{各面在同一顶点的切向扇区。内部顶点铺满一周，光滑边界顶点铺满半周，原有角点只铺满原内角。图中的圆弧仅用于标示扇区，不是曲面的分解边。}
	\label{g1:fig:gb-vertex-angles}
\end{figure}


为统一记号，给每个边界顶点置 \(\alpha_v=\pi-\beta_v\)，在新增的光滑分点处置 \(\alpha_v=0\)。则
\begin{equation}\label{g1:eq:gb-interior-angle-count}
	\sum_{j,v}\beta_{j,v}
	=2\pi V_{\mathrm i}+\pi V_{\mathrm b}
	-\sum_{v\in\partial M}\alpha_v.
\end{equation}
这里的角度是度量角，选择正交坐标后才与图中的 Euclid 角一致；角能相加是同一个切平面中的事实，不要求坐标图把曲面等距铺平。

\subsection{Euler 示性数的出现}
\label{g1:subsec:0904:03}

每个闭圆盘的边数等于角的出现次数。内部边出现两次，边界边出现一次，故总出现次数为 \(2E_{\mathrm i}+E_{\mathrm b}\)。把每个面角换成外角 \(\pi-\beta_{j,v}\)，再代入\eqref{g1:eq:gb-interior-angle-count}，得到
\begin{align*}
	\sum_{j,v}(\pi-\beta_{j,v})
	&=\pi(2E_{\mathrm i}+E_{\mathrm b})-2\pi V_{\mathrm i}-\pi V_{\mathrm b}
	+\sum_{v\in\partial M}\alpha_v\\
	&=2\pi(E_{\mathrm i}-V_{\mathrm i})+\sum_{v\in\partial M}\alpha_v.
\end{align*}
第二步用到 \(E_{\mathrm b}=V_{\mathrm b}\)：\cref{g1:prop:boundary-circles}说明紧曲面的边界是有限个圆周，每条圆周分成多少条边，就有多少个分点。这也解释了为何边界顶点不能与内部顶点一起按 \(2\pi\) 计角。

\ProseParagraphBreak

现在对每块 \(D_j\) 使用\cref{g1:thm:local-gauss-bonnet}。分解引理保证它是落在单张图中的闭圆盘，边界具有所要求的正则性和非零角，因此每块的右端都是 \(2\pi\)。对各块相加，面积积分可加、公共边按\eqref{g1:eq:gb-edge-cancellation}抵消，故
\[
	\int_MK\,\mathrm{d}A+\int_{\partial M}k_g\,\mathrm{d}s
	+\sum_{v\in\partial M}\alpha_v
	=2\pi(F-E_{\mathrm i}+V_{\mathrm i}).
\]
由于 \(E_{\mathrm b}=V_{\mathrm b}\)，右端括号等于
\[
	F-(E_{\mathrm i}+E_{\mathrm b})+(V_{\mathrm i}+V_{\mathrm b})=V-E+F=\chi(M).
\]
最后一步正是\cref{g1:prop:euler-triangulation-independent}：它不仅证明三角剖分的交错和不变，也证明任何有限闭盘分解的交错和都等于同一个 Euler 示性数。这里已经构造出真正的闭盘分解，因而满足该命题的条件。

\ProseParagraphBreak

可以把这段计数理解为两笔账。曲率积分记下各小片内部与边缘的几何；面片外角的总和则记下这些小片在顶点附近如何拼接。把人为切出的边和角消去后，几何账中留下总曲率，组合账中留下 \(\chi(M)\)。

\subsection{闭曲面与带边界曲面的统一公式}
\label{g1:subsec:0904:04}

\begin{theorem}[Gauss--Bonnet]\label{g1:thm:global-gauss-bonnet}
	设 \(M\) 为紧、Hausdorff、第二可数的光滑度量曲面，允许不连通、不可定向以及分片 \(C^2\) 边界。边界由有限条正则弧组成，角点的内角 \(\beta_v\) 属于 \((0,2\pi)\)。度量在边界及角点的局部光滑环境中光滑延拓。令 \(K\) 为 Gauss 曲率，\(\mathrm{d}A\) 为度量面积测度，沿每条边界弧取指向 \(M\) 内侧的单位切向法向 \(\nu\)，并定义
	\[
		k_g=\langle\nabla_TT,\nu\rangle,\qquad \alpha_v=\pi-\beta_v.
	\]
	这里 \(T\) 是沿该弧任一方向的单位切向量。则
	\begin{equation}\label{g1:eq:global-gauss-bonnet}
		\int_MK\,\mathrm{d}A+\int_{\partial M}k_g\,\mathrm{d}s+\sum_v\alpha_v
		=2\pi\chi(M).
	\end{equation}
	当 \(M\) 可定向时，边界按区域在左的方向行进，上述 \(k_g\) 就是 \(\langle\nabla_TT,JT\rangle\)。不连通时，\(\chi(M)\) 为各连通分支的 Euler 示性数之和。
\end{theorem}
\begin{proof}
	对每个连通分支使用\cref{g1:lem:smooth-coordinate-disk-decomposition}，按前两小节的边项抵消和顶点计数相加，即得\eqref{g1:eq:global-gauss-bonnet}。局部面积测度与 \(K\) 不因反转定向而改变；内侧测地曲率的抵消已经逐边证明，故不需要全局定向。最后，紧曲面只有有限个连通分支，这是\cref{g1:prop:boundary-circles}的另一项结论；将有限个公式相加即可。
\end{proof}

若边界光滑，角项为空；若 \(M\) 是闭曲面，边界积分也为空，公式成为
\begin{equation}\label{g1:eq:closed-gauss-bonnet}
	\int_MK\,\mathrm{d}A=2\pi\chi(M).
\end{equation}
另一方面，若 \(M\) 本身是一个圆盘，即使它不落在任何一张坐标图内，也有 \(\chi(M)=1\)，所以仍得到右端为 \(2\pi\) 的带边界公式。单图条件是局部证明所用的条件，经粘合后便不再是整体结论的限制。

\begin{example}\label{g1:exm:flat-mobius-gauss-bonnet}
	取带 Euclid 度量的矩形 \([0,L]\times[-a,a]\)，其中 \(L,a>0\)，按 \((0,y)\sim(L,-y)\) 识别两条竖边。因为识别映射在邻近条带上由刚性变换 \((x,y)\mapsto(x+L,-y)\) 给出，度量光滑粘合成一条 Möbius 带。这里的度量由矩形粘合定义，不需要把它作为 \(\mathbb R^3\) 中的曲面实现。

	局部图都是 Euclid 图，所以 \(K=0\)。上下两条水平边接成同一条光滑边界，各局部段都是直线，故 \(k_g=0\)，识别处也没有角点。\cref{g1:prop:surface-euler-formulas}给出 \(\chi(N_{1,1})=2-1-1=0\)，于是公式两边都是零。这个例子中不存在连续的全局定向，但“指向带内侧”的边界法向仍然存在，正好对应定理中采用的表述。
\end{example}

有边界时不能只看面积曲率。平面圆盘的 \(K\) 处处为零而 \(\chi=1\)，缺少的 \(2\pi\) 全由边界转动提供；平面环带的内圈则带来负的边界测地曲率，抵消外圈的贡献。角项承担同样的补偿作用，所以把光滑边界换成多边形边界后，公式的右端仍是同一个拓扑数。
